\subsection*{静电平衡、空腔与屏蔽}

导体内有可自由移动的电荷。若内部存在非零电场，自由电荷会继续运动，故静电平衡时导体内部
$\vec E=0$，整个导体及其表面为等势体。表面切向电场必须为零，否则表面电荷也会运动；电场只能
沿法线。用跨过表面的薄柱高斯面可得外侧法向场
\[
 E_n=\frac\sigma{\varepsilon_0}.
\]
多余净电荷分布在外表面，但分布通常不均匀：几何越尖锐处，为维持等势所需的面电荷密度越大，
邻近场强也越强。

导体空腔内若没有自由电荷，腔壁又是同一等势面，唯一性定理保证腔内电势为常量、场为零，外界
电场不能穿入，这就是静电屏蔽。若腔内放入电荷 $q$，腔壁总感应电荷为 $-q$；导体原有净电荷与
$q$ 的影响共同决定外表面总电荷。腔内场受腔内电荷和腔形决定，而在封闭导体中移动腔内电荷，
不会改变由外部边界条件唯一确定的外部场。

唯一性定理的内容是：给定区域内电荷分布，并在边界上给定电势或法向电场后，满足 Poisson 或
Laplace 方程的解至多一个。证明可设有两个解，考察其差；差函数满足 Laplace 方程且边界为零，
由能量积分或极值原理只能处处为零。此定理让我们不必直接解偏微分方程：只要构造一个满足场方程
和全部边界条件的候选解，它就是物理解。

\subsection*{电像法与受力、能量}

电像法用导体区域以外的虚拟电荷替代未知表面电荷，使导体边界成为规定等势面。接地无限平面前
的点电荷 $q$，可在镜像位置放置 $-q$；真实半空间中的电势与场便完全正确。像电荷不是真实电荷，
不能在真实空间里另算其电荷守恒，也不能求“像电荷受到的力”。真实电荷所受力可由像电荷产生的
场计算。

能量计算要特别注意装配过程。点电荷靠近平面时，像电荷的位置会随真实电荷移动，系统能量不是
简单的真实电荷与像电荷相互作用能，而是其一半：
\[
 U=\frac12q\varphi_{\rm image}(\vec r_q).
\]
对固定位置直接求力则不需要这个 $1/2$。球面边界的像电荷位置和电量由反演关系决定；接地球与
孤立中性球的边界条件不同，后者还要增加中心像电荷以保证导体总电荷为零。

等离子体中电子相对几乎不动的离子背景发生微小整体位移时，会出现把电子拉回平衡位置的电场。
电子密度为 $n$ 时，线性化运动方程给出等离子体频率
\[
 \omega_p=\sqrt{\frac{ne^2}{m_e\varepsilon_0}}.
\]
该频率与振幅无关是小振动近似的结果；它说明导体屏蔽并非“电荷瞬时到位”，自由载流子有自身的
集体响应时间尺度。

